\documentclass[11pt,a4paper]{article}

% ---------------------------------------------------------------------------
% Tiếng Việt (biên dịch bằng pdfLaTeX, ví dụ trên Overleaf)
% ---------------------------------------------------------------------------
\usepackage[T5]{fontenc}
\usepackage[utf8]{inputenc}
\usepackage[vietnamese]{babel}

\usepackage[a4paper,top=2.2cm,bottom=2.6cm,left=2.6cm,right=2.6cm,footskip=1.2cm]{geometry}
\usepackage{graphicx}
\usepackage{xcolor}
\usepackage{tikz}
\usetikzlibrary{shadows}
\usepackage{eso-pic}
\usepackage{fancyhdr}
\usepackage{enumitem}
\usepackage{upquote}
\usepackage{array}
\usepackage{parskip}

\definecolor{aivngreen}{HTML}{2F4F1F}
\definecolor{aivngray}{HTML}{6B6B6B}

% Khung viền quanh trang
\AddToShipoutPictureBG{%
  \begin{tikzpicture}[remember picture,overlay]
    \draw[line width=0.8pt]
      ([xshift=0.35cm,yshift=-0.35cm]current page.north west)
      rectangle
      ([xshift=-0.35cm,yshift=0.35cm]current page.south east);
  \end{tikzpicture}%
}

% Chân trang
\pagestyle{fancy}
\fancyhf{}
\renewcommand{\headrulewidth}{0pt}
\renewcommand{\footrulewidth}{0.8pt}
\fancyfoot[L]{\textbf{SĐT/Zalo: 0911118758 (Dr. Hồng Phúc)}}
\fancyfoot[R]{\textbf{Trang \thepage}}
\fancypagestyle{plain}{%
  \fancyhf{}
  \renewcommand{\headrulewidth}{0pt}
  \renewcommand{\footrulewidth}{0.8pt}
  \fancyfoot[L]{\textbf{SĐT/Zalo: 0911118758 (Dr. Hồng Phúc)}}
  \fancyfoot[R]{\textbf{Trang \thepage}}
}

% Logo: dùng file logo-aivn.png nếu có, nếu không vẽ tạm một vòng tròn
\newcommand{\aivnlogo}{%
  \IfFileExists{logo-aivn.png}{%
    \includegraphics[height=1.6cm]{logo-aivn.png}%
  }{%
    \begin{tikzpicture}
      \node[circle,fill=white,draw=black!15,drop shadow,minimum size=1.5cm,
            font=\fontsize{24}{24}\selectfont\bfseries\sffamily,
            text=aivngreen] {AI};
    \end{tikzpicture}%
  }%
}

% Câu hỏi và lựa chọn
\newcounter{question}
\newcommand{\question}{\par\medskip\stepcounter{question}\noindent\textbf{Câu \thequestion.}\ }
\newenvironment{choices}
  {\begin{enumerate}[label=\alph*),leftmargin=2.2em,itemsep=1pt,topsep=3pt]}
  {\end{enumerate}}
\newcommand{\code}[1]{\texttt{#1}}

\begin{document}
\thispagestyle{plain}

% ---------------------------------------------------------------------------
% Phần đầu trang
% ---------------------------------------------------------------------------
\noindent
\begin{minipage}[c]{1.9cm}\aivnlogo\end{minipage}%
\begin{minipage}[c]{8cm}
  {\bfseries AI VIET NAM}\\[1pt]
  {\small\color{aivngray}@aivietnam.edu.vn}
\end{minipage}

\vspace{0.9cm}
\noindent\rule{\linewidth}{0.6pt}
\begin{center}
  {\LARGE\bfseries Quiz: Transformation with dbt}
\end{center}
\noindent\rule{\linewidth}{0.6pt}

\vspace{0.4cm}
\begin{center}
  \textbf{Lê Thanh Bảo Trân}\hspace{2.2cm}\textbf{Huỳnh Trần Anh Thư}\\[6pt]
  \textbf{Proofreading:}
\end{center}
\vspace{0.6cm}

\section*{Hướng dẫn}

Quiz gồm 12 câu, bám theo năm phần của bài \textit{Transformation with dbt}: vì sao có
dbt, dbt hoạt động thế nào, thiết kế dữ liệu, dữ liệu đáng tin cậy và các tính năng
nâng cao. Câu trắc nghiệm chọn một đáp án. Câu ghép nối ghi theo dạng A--1. Câu tự
luận trả lời ngắn trong một đến hai câu. Đáp án và giải thích ở cuối tài liệu.

% ---------------------------------------------------------------------------
\section*{Phần I--II. dbt là gì, hoạt động thế nào}
% ---------------------------------------------------------------------------

\question Vì sao các tổ chức chuyển từ ETL sang ELT?
\begin{choices}
  \item SQL chạy nhanh hơn Python.
  \item Lưu trữ và tính toán trên cloud warehouse đã rẻ và mạnh, nên có thể nạp dữ
        liệu thô vào trước rồi biến đổi ngay trong warehouse.
  \item ETL không xử lý được dữ liệu dạng CSV.
  \item dbt bắt buộc phải dùng ELT.
\end{choices}

\question Phát biểu nào sau đây \emph{sai} về dbt?
\begin{choices}
  \item dbt biên dịch SQL kết hợp Jinja thành SQL thuần.
  \item dbt tự xử lý dữ liệu bằng engine riêng của nó.
  \item dbt điều phối thứ tự chạy model theo DAG.
  \item dbt không thực hiện bước extract và load.
\end{choices}

\question (Tự luận) Một model trong dbt là gì? Vì sao trong file model không cần viết
\code{CREATE TABLE}?

\question (Tự luận) Hàm \code{source()} và \code{ref()} khác nhau như thế nào? Hàm nào
là điểm bắt đầu của DAG?

\question Ghép mỗi materialization với mô tả phù hợp.

\begin{center}
\begin{tabular}{>{\bfseries}l l @{\hspace{1.2cm}} >{\bfseries}l p{8cm}}
  A. & view        & 1. & Không tạo object, được chèn vào model khác dưới dạng CTE. \\
  B. & table       & 2. & Chỉ lưu câu query, mỗi lần đọc thì tính lại. \\
  C. & incremental & 3. & Lần đầu tạo bảng, các lần sau chỉ xử lý dòng mới hoặc dòng thay đổi. \\
  D. & ephemeral   & 4. & Chạy query, lưu kết quả và build lại toàn bộ mỗi lần. \\
\end{tabular}
\end{center}

\question (Tự luận) Một model incremental lọc theo điều kiện
\code{updated\_at > max(updated\_at)} và không khai báo \code{unique\_key}. Nêu hai lỗi
có thể xảy ra và cách phòng cho từng lỗi.

% ---------------------------------------------------------------------------
\section*{Phần III. Thiết kế dữ liệu}
% ---------------------------------------------------------------------------

\question Việc nào \emph{không} nên làm ở tầng staging?
\begin{choices}
  \item Đổi tên cột, ví dụ \code{id} thành \code{order\_id}.
  \item Ép kiểu dữ liệu, ví dụ chuỗi thành \code{date}.
  \item Join với bảng khác hoặc dùng \code{group by}.
  \item Chuẩn hoá giá trị bằng \code{lower()}, \code{trim()}.
\end{choices}

\question (Tự luận) Nêu ba trường hợp nên tạo model intermediate thay vì viết logic
thẳng trong mart.

\question (Tự luận) Trong star schema, fact table và dimension table khác nhau ở điểm
nào?

\question (Tự luận) Grain là gì? Fan-out xảy ra khi nào, và vì sao lỗi này khó phát
hiện?

\question Đơn hàng A có 2 sản phẩm và 2 lần thanh toán, mỗi lần 100. Join thẳng bảng
\code{order\_items} với bảng \code{payments} theo \code{order\_id}. Mỗi cách tính dưới
đây cho kết quả bao nhiêu, và cách nào đúng?
\begin{choices}
  \item \code{sum(amount)} trên kết quả join.
  \item \code{sum(distinct amount)} trên kết quả join.
  \item Gom \code{payments} về grain đơn hàng (\code{group by order\_id}) trước, rồi mới
        join.
\end{choices}

% ---------------------------------------------------------------------------
\section*{Phần IV--V. Kiểm thử và tính năng nâng cao}
% ---------------------------------------------------------------------------

\question (Tự luận) Unit test, data test và model contract mỗi loại kiểm tra điều gì?
Vì sao không loại nào thay thế được loại nào?

% ===========================================================================
\newpage
\section*{Đáp án và giải thích}
% ===========================================================================

\begin{enumerate}[label=\textbf{Câu \arabic*.},leftmargin=*,itemsep=6pt]

  \item \textbf{b.} Trong ETL, dữ liệu được biến đổi trên máy chủ trung gian vì lưu
  trữ và tính toán trong warehouse on-premise đắt. Cloud warehouse và các engine cột
  như DuckDB làm cho lưu trữ rẻ, tính toán mạnh và tách rời nhau. Dữ liệu thô luôn còn
  trong kho, nên khi logic sai chỉ cần chạy lại phần biến đổi.

  \item \textbf{b.} dbt không xử lý dữ liệu. Mọi phép tính do warehouse thực hiện trên
  câu SQL mà dbt sinh ra, nên hiệu năng phụ thuộc vào warehouse.

  \item Model là một file \code{.sql} chứa đúng một câu \code{SELECT}, tên file là tên
  bảng hoặc view được tạo. dbt tự bọc câu \code{SELECT} thành lệnh tạo object phù hợp
  (\code{create view}, \code{create table}\ldots) theo materialization đã cấu hình.

  \item \code{source()} trỏ tới dữ liệu thô được khai báo trong file YAML. Đây là điểm
  bắt đầu của DAG và cho phép kiểm tra độ tươi bằng \code{dbt source freshness}.
  \code{ref()} trỏ tới một model khác. Mỗi lần gọi \code{ref()} là ngầm khai báo một
  phụ thuộc, và tên model được đổi thành tên đầy đủ theo target đang chạy (dev hoặc
  prod).

  \item \textbf{A--2, B--4, C--3, D--1.}

  \item Lỗi thứ nhất là dữ liệu đến trễ bị bỏ qua vĩnh viễn: dòng có \code{updated\_at}
  nhỏ hơn mốc hiện tại nhưng được nạp muộn sẽ không bao giờ thoả điều kiện lọc. Cách
  phòng là lọc lùi một khoảng an toàn, ví dụ \code{max(updated\_at) - interval 3 day}.
  Lỗi thứ hai là dữ liệu bị chèn trùng: không có \code{unique\_key} thì dbt chỉ thêm
  dòng, nên một đơn đổi trạng thái xuất hiện hai lần. Cách phòng là khai báo
  \code{unique\_key} để dbt merge. Hai cách phải đi cùng nhau, vì lọc lùi nghĩa là
  chủ động lấy lại dòng cũ. Khi bảng đã sai, chạy \code{dbt build --full-refresh}.

  \item \textbf{c.} Join làm staging phụ thuộc vào bảng khác, còn \code{group by} làm
  đổi grain. Staging chỉ làm sạch từng bảng nguồn theo quan hệ 1:1.

  \item (1) Logic quá dài làm model mart khó đọc. (2) Cùng một logic được dùng lại ở
  nhiều mart. (3) Cần đổi grain trước khi join.

  \item Fact table ghi lại sự kiện đo đếm được (một lần bán, một lần thanh toán), chứa
  các chỉ số như số lượng, số tiền và khoá ngoại tới dimension. Dimension table mô tả
  ngữ cảnh của sự kiện (ai mua, mua gì, khi nào), gồm các cột dùng để lọc và nhóm.

  \item Grain là câu trả lời cho câu hỏi ``một dòng trong bảng này đại diện cho cái
  gì''. Fan-out xảy ra khi join hai quan hệ một-nhiều theo cùng một khoá, làm số dòng
  bị nhân lên và mọi phép \code{sum()} phía sau bị đội lên. Lỗi khó phát hiện vì câu
  SQL vẫn chạy không báo lỗi; phần lớn bản ghi có hệ số nhân bằng 1 nên tổng chỉ lệch
  vài phần trăm; và lỗi có thể xuất hiện muộn khi dữ liệu nguồn thay đổi (ví dụ có
  thêm tính năng trả góp) dù model không đổi.

  \item Kết quả join có $2 \times 2 = 4$ dòng.
  (a) \code{sum(amount)} $= 400$, sai vì mỗi lần thanh toán bị lặp theo số sản phẩm.
  (b) \code{sum(distinct amount)} $= 100$, sai vì hai lần thanh toán có cùng số tiền
  nên bị gộp làm một.
  (c) $= 200$, đúng. Gom về cùng grain rồi mới join là cách làm chuẩn.

  \item Unit test kiểm tra logic SQL trên vài dòng dữ liệu giả, bắt sai công thức và
  trường hợp biên. Model contract kiểm tra hình dạng (tên cột, kiểu dữ liệu) ngay lúc
  build, sai thì dbt dừng trước khi tạo bảng. Data test kiểm tra dữ liệu thật sau khi
  build, bắt dữ liệu nguồn bẩn, fan-out, khoá ngoại mồ côi. Mỗi loại bắt một nhóm lỗi
  khác nhau: logic đúng vẫn có thể gặp dữ liệu bẩn, và dữ liệu sạch vẫn có thể bị tính
  sai công thức.

\end{enumerate}

\end{document}
